 Classification: Biological Sciences


Keywords (Not sure what you guys use as keywords): 
phenology, climate change, statistical methodology, stopping times, variance estimation
% emw4July -- these seem fine; you're supposed to just try to pick words that would help old school search engines, so things your paper is about that are not in the title or abstract (ideally). 

Significance statement (120 words max, I treated this as a ELI5 exercise—not sure if appropriate. I think I used the text you from the draft.):

% emw4July -- the  ELI5 exercise angle is great! But it read a little theoretical so I tried to get a more concrete result in at the end. 
Springtime events like flowering depend on temperature, with warmer temperatures shifting events earlier. Climate change has thus advanced spring events. But the rate of advance has slowed in recent years, alongside increasing variance. While research has suggested complex factors to explain these trends, we show that they follow a simple thermal-sum model in which an event occurs once cumulative temperature reach a threshold. When average temperatures increase day over day, as they often do near the spring equinox, events are usually tightly synchronized. But when warming shifts those events closer to the winter solstice, event times become more variable. Our results suggest this regime shift to a `scattered spring' may increase by the end of the century, reducing ecological synchrony. 

Competing interest statement: The authors declare no competing interests. 

Data Sharing Plans: All code and data are documented in the author's GitHub repository, which is readily available upon request, and will be made pubicly available via Zenodo, KNB or similar upon publication. % emw4July -- slight teaks, I don't care what repo we use as long as it not FigShare or something like that. 

Funding information: The authors declare no sources of funding.

Reviewers:

For biology, ...

For statistical methodology, we could suggest Michael Baron (American University) and Yaakov Malinovsky (University of Maryland) who recently hosted the International Workshop in Sequential Methodologies. The arguments we make in our paper build on the sequential methodologies literature, and we believe these reviewers would have a good perspective on appropriate methods and practices. % emw4July -- seems good. 

Letter (again, no idea what to stress to a biologist):

Dear Editor,

% emw4July -- less passive voice for biologists, more repeating the results and make it more exciting! I always say you go BIG in the cover letter, bigger than you go in the article. You're trying to get the article OUT for review here. See what you think, I also usually cite a couple recent papers in PNAS or higher to convince the editor the work is of interest to high profile journals.
In this manuscript, we address a central puzzle in plant biology and phenology that has become increasingly apparent with climate change. Scientists have long recognized that warmer temperatures drive earlier springtime events, such as flowering and insect emergence. In recent years, however, two additional patterns have emerged: the rate at which these events advance has slowed, and their timing has often become more variable. % Without accounting for both patterns, it is difficult to predict how ecosystems will respond as climates continue to warm.

Researchers have invoked complex models to explain these patterns. An extensive and rapidly growing body of work has added interacting factors to a basic thermal sum model, in which biological events are triggered once cumulative temperatures reach a threshold (this includes the common growing degree day mdoel). While over a decade of work has debated whether photoperiod, winter temperatures, or cold hardiness may underlie these patterns \citep{fu2015declining,kovaleskipreprint}, no work has tested the full predictions for shifting event times from the simple thermal sum model. 

Here, working across biology and statistics, we show that these patterns are consistent with the simple thermal-sum model. Our model predicts both slower advances and increased variance with increased warming, and we validate these model predictions with experimental data and observational lilac flowering data over the last 70 years across the US. In contrast to predictions from current models of biological events, our framework implies that as springtime events advance with anthropogenic warming, they enter a new temperature regime in which the typical time between events in the same ecosystem can increase from a few days to a month or more. % Importantly, our results reveal the key mechanism that determines whether variation declines: the temperature regime under which accumulation occurs. When average temperatures increase day over day, as they often do near the equinox, springtime events tend to be tightly synchronized. But when warming shifts those events closer to the solstice, where average temperatures are comparatively stationary, event times become more variable. 

These results challenge research that has assumed interactive factors beyond temperature explain these trends. Further, they imply that commonly used statistical tools that assume linear responses or constant variance may be misspecified and can lead to biased or inaccurate forecasts. Our findings thus challenge current approaches. Yet, at the same time, they can lead to major advances in forecasting: we provide a simple and transparent way to encode the thermal-sum model, on top of which complexity may be added and better studied.

The work is interdisciplinary. It builds on biological theory from phenology and on mathematical results from the study of stopped random walks. At the same time, our goal is to reach a broad audience working at the intersection of ecology, agriculture, plant biology, climate change, public policy, and statistical methodology. We have therefore simplified the presentation while retaining the key theoretical insight: the same temperature-threshold mechanism that has long been used to explain the timing of spring events also predicts when those events will become less sensitive to warming and less synchronized across space.

We believe the manuscript will be of interest to readers concerned with fundamental plant biology, environmental and ecological forecasting, the impact of climate change, and the statistical modeling of biological event times.

Thank you for your consideration,

Jonathan Auerbach, Lizzie Wolkovich (not sure how you want to sign here), and Andrew Gelman
